\chapter{Overview}
\label{chap:overview}

% archon:covers MR4513142ThereIsNoEnriquesSurfaceOverTheIntegers.lean MR4513142ThereIsNoEnriquesSurfaceOverTheIntegers/Basic.lean

% Chapter-local macros (plan agent: add the following to common.tex as shared macros)
% \ZZ, \FF, \RR, \QQ, \Spec, \Num, \Br, \foY, \foX, \shM, \Enr, \ra, \lra, \quadand, \bY, \GG
\providecommand{\ZZ}{\mathbb{Z}}
\providecommand{\FF}{\mathbb{F}}
\providecommand{\RR}{\mathbb{R}}
\providecommand{\QQ}{\mathbb{Q}}
\providecommand{\Spec}{\mathrm{Spec}}
\providecommand{\Num}{\mathrm{Num}}
\providecommand{\Br}{\operatorname{Br}}
\providecommand{\GG}{\mathbb{G}}
\providecommand{\foY}{\mathfrak{Y}}
\providecommand{\foX}{\mathfrak{X}}
\providecommand{\shM}{\mathscr{M}}
\providecommand{\Enr}{\mathrm{Enr}}
\providecommand{\ra}{\rightarrow}
\providecommand{\lra}{\longrightarrow}
\providecommand{\quadand}{\quad\text{and}\quad}
\providecommand{\bY}{\bar{Y}}

This chapter covers the formalization of the main result of
S.~Schr\"oer, \textit{There is no Enriques surface over the integers},
Amer.\ J.\ Math.\ 144 (2022), 1585--1677 (MR4513142, arXiv:2004.07025).
The paper establishes that the moduli stack \(\shM_{\Enr}\) of Enriques surfaces has empty fiber category
\(\shM_{\Enr}(\ZZ) = \varnothing\): there is no flat proper family of Enriques surfaces over \(\Spec(\ZZ)\).
The proof is long and indirect.
It reduces the problem to showing that no Enriques surface over the prime field \(\FF_2\) is
simultaneously non-exceptional and has constant Picard scheme,
a result established by classifying genus-one fibrations and their fiber configurations.


%-----------------------------------------------------------------
\section{The local system of numerical classes}
%-----------------------------------------------------------------

\begin{definition}[Numerically trivial Picard scheme]
  \label{def:pic-tau}
  \lean{MR4513142ThereIsNoEnriquesSurfaceOverTheIntegers.PicTau}
  \uses{def:picard-scheme}
  % SOURCE: [Schröer, arXiv:2004.07025], §1, lines 593–596 (read from references/MR4513142-enriques-integers.tex)
  % SOURCE QUOTE: "Let $X$ be a proper scheme,   $\Pic_{X/k}$ the
  % Picard scheme, and $\Pic^\tau_{X/k}$ be the open subgroup scheme parameterizing numerically trivial
  % invertible sheaves  (see for example \cite{SGA 6}, Expos\'e XII, Corollary 1.5 and Expos\'e XIII, Theorem 4.7, compare also
  % \cite{Laurent; Schroeer 2021}, Section 2)."
  \textit{Source: Schr\"oer, §1 (citing SGA 6, XII.1.5 and XIII.4.7).}
  The \emph{numerically trivial Picard scheme} \(\Pic^{\tau}_{\foY/R}\) is the open subgroup scheme of
  \(\Pic_{\foY/R}\) parameterizing numerically trivial invertible sheaves.
  For a family of Enriques surfaces, its structure morphism over \(\Spec(R)\) is locally free of rank two.
\end{definition}

\begin{definition}[Local system of numerical classes]
  \label{def:num-local-system}
  \lean{MR4513142ThereIsNoEnriquesSurfaceOverTheIntegers.NumLocalSystem}
  \uses{def:picard-scheme, def:pic-tau}
  % SOURCE: [Schröer, arXiv:2004.07025], §1, lines 598–610 (read from references/MR4513142-enriques-integers.tex)
  % SOURCE QUOTE: "the resulting quotient
  % $$
  % \Num_{X/k}=(\Pic_{X/k})/(\Pic^\tau_{X/k})
  % $$
  % is an \'etale group scheme.
  % As such, it corresponds to the Galois representation on the stalk
  % $$
  % \Num_{X/k}(k^\alg)=\Num_{X/k}(k^\sep)=\ZZ^{\oplus\rho},
  % $$
  % where  $k^\sep\subset k^\alg$ are chosen separable and algebraic closures,
  % and  $\rho\geq 0$ denotes the Picard number for the base-change $\bX=X\otimes k^\alg$.
  % ...
  % One may call  $\Num_{X/k}$  the \emph{local system  of numerical classes}."
  \textit{Source: Schr\"oer, §1.}
  The \emph{local system of numerical classes} is the \'etale group scheme
  \[
    \Num_{\foY/R} \;=\; \Pic_{\foY/R} \,/\, \Pic^{\tau}_{\foY/R}.
  \]
  It corresponds to the Galois representation on the stalk
  \(\Num_{\foY/R}(k^{\mathrm{alg}}) = \ZZ^{\oplus \rho}\), where \(\rho\) is the geometric Picard number.
  For a family of Enriques surfaces, \(\Num_{\foY/R}\) is a local system of free abelian groups of rank \(\rho = 10\).
\end{definition}

\begin{theorem}[Descent of projective contractions]
  \label{thm:projective-contractions}
  \lean{MR4513142ThereIsNoEnriquesSurfaceOverTheIntegers.projectiveContractions}
  \uses{def:num-local-system}
  % SOURCE: [Schröer, arXiv:2004.07025], §1, lines 714–718 (read from references/MR4513142-enriques-integers.tex)
  % SOURCE QUOTE: "\begin{theorem}
  % \mylabel{projective contractions}
  % Suppose  the local system $\Num_{X/k}$ is constant.
  % Then every contraction $\bX\ra \bY$  to some projective scheme $\bY$ is isomorphic to the base-change of a contraction $X\ra Y$.
  % \end{theorem}"
  \textit{Source: Schr\"oer, §1, Thm.~1.3.}
  Let \(X\) be a proper scheme over a field \(k\) such that the local system of numerical classes
  \(\Num_{X/k}\) is constant. Then every contraction \(\bX \to \bY\) to a projective scheme \(\bY\)
  (over the algebraic closure \(\bX = X \otimes_k k^{\mathrm{alg}}\)) is isomorphic to the base-change of a
  contraction \(X \to Y\) defined over \(k\). Here a \emph{contraction} of \(X\) is a proper scheme \(Z\) with a
  morphism \(g : X \to Z\) satisfying \(\cO_Z = g_*(\cO_X)\).
\end{theorem}
\begin{proof}
  Each invertible sheaf \(\shL\) on \(X\) yields the graded ring \(R(X, \shL) = \bigoplus_{t \geq 0} H^0(X, \shL^{\otimes t})\),
  and if \(\shL\) is semiample then \(\Proj R(X, \shL)\) is a projective contraction, every projective contraction
  arising this way (with \(\shL\) unique up to preimages of semiample sheaves). By \cite{EGA IVc}, Theorem~8.8.2,
  the contraction \(\bX \to \bY\) descends to some finite extension \(k \subset k'\); enlarging \(k'\) to a splitting
  field and treating the Galois and purely inseparable cases separately, constancy of \(\Num_{X/k}\) lets the defining
  semiample sheaf descend to \(X\), giving the desired contraction over \(k\).
\end{proof}

%-----------------------------------------------------------------
\section{Enriques surfaces and exceptional Enriques surfaces}
%-----------------------------------------------------------------

\begin{definition}[Enriques surface over a field]
  \label{def:enriques-surface}
  \lean{MR4513142ThereIsNoEnriquesSurfaceOverTheIntegers.IsEnriquesSurface}
  \uses{def:pic-tau, def:num-local-system}
  % SOURCE: [Schröer, arXiv:2004.07025], §4, lines 1052–1062 (read from references/MR4513142-enriques-integers.tex)
  % SOURCE QUOTE: "Throughout, we fix an algebraically closed ground field $k$ of characteristic $p\geq 0$.
  % An \emph{Enriques surface}  is a smooth proper  scheme $Y$ with
  % $$
  % h^0(\O_Y)=1\quadand  c_1=0\quadand b_2=10.
  % $$
  % The first condition means that $Y$ is connected, and  the second condition
  %  signifies that the dualizing sheaf $\omega_Y$ is numerically trivial.
  % In the Enriques classification of algebraic surfaces, Bombieri and Mumford  showed in  \cite{Bombieri; Mumford 1976}, \S 3    that
  % the   Picard scheme $\Pic^\tau_{Y/k}$ of numerically trivial sheaves is a finite group scheme
  % of order two, whose group of rational points is generated by the dualizing sheaf $\omega_Y$.
  % Moreover, $\Num(Y)$ is a  free abelian group of rank $\rho=10$, the Betti numbers satisfies $b_1=b_3=0$ and $b_2=\rho$."
  \textit{Source: Schr\"oer, §4 (citing Bombieri--Mumford 1976).}
  An \emph{Enriques surface} over an algebraically closed field \(k\) is a smooth proper scheme \(Y\) satisfying
  \[
    h^0(\cO_Y) = 1, \quad c_1 = 0, \quad b_2 = 10.
  \]
  Here \(c_1 = 0\) means the dualizing sheaf \(\omega_Y\) is numerically trivial.
  Equivalently (Bombieri--Mumford), \(\Pic^{\tau}_{Y/k}\) is a finite group scheme of order two generated by \(\omega_Y\),
  and \(\Num(Y) \cong \ZZ^{\oplus 10}\) with Betti numbers \(b_1 = b_3 = 0\).
  Over algebraically closed fields of characteristic \(\neq 2\), this is equivalent to
  \(\omega_Y \not\cong \cO_Y\) and \(\omega_Y^{\otimes 2} \cong \cO_Y\).
\end{definition}

\begin{definition}[Exceptional Enriques surface]
  \label{def:exceptional-enriques}
  \lean{MR4513142ThereIsNoEnriquesSurfaceOverTheIntegers.IsExceptionalEnriquesSurface}
  \uses{def:enriques-surface, def:pic-tau}
  % SOURCE: [Schröer, arXiv:2004.07025], §4, lines 1130–1147 (read from references/MR4513142-enriques-integers.tex)
  % SOURCE QUOTE: "In case that $P=\Pic^\tau_{Y/k}$ is unipotent, in other words, the Cartier dual $G=\uHom(P,\GG_m)$ is local,
  % one says that $Y$ is \emph{simply-connected}."
  % SOURCE QUOTE: "The simply-connected Enriques surface can be divided into two   subclasses depending on
  % properties of the normalization $\nu:X'\ra X$. The induced projection $X'\ra Y$ is purely inseparable and flat of degree two,
  % hence the ramification locus for $\nu$ is an effective Cartier divisor $R'\subset X'$, giving $\omega_{X'}=\O_{X'}(-R')$.
  % Ekedahl and Shepherd-Barron \cite{Ekedahl; Shepherd-Barron 2004}
  % observed that $R'$ is the preimage of some effective Cartier divisor $C\subset Y$
  % referred to as the  \emph{conductrix}, and  $2C$ is called the \emph{biconductrix}.
  % They  called an Enriques surface \emph{exceptional} if it is simply-connected and
  % the biconductrix has $h^1(\O_{2C})\neq 0$."
  % SOURCE QUOTE (table, lines 1111–1127): in characteristic $p=2$ the three types are $P=\mu_2$ (ordinary,
  % $\pi_1=\ZZ/2\ZZ$), $P=\ZZ/2\ZZ$ (classical, $\pi_1$ trivial) and $P=\alpha_2$ (supersingular, $\pi_1$ trivial).
  \textit{Source: Schr\"oer, §4 (citing Ekedahl--Shepherd-Barron 2004).}
  In characteristic \(2\), an Enriques surface \(Y\) over an algebraically closed field is \emph{simply-connected}
  if its torsion Picard scheme \(\Pic^{\tau}_{Y/k}\) is unipotent, equivalently its Cartier dual is local; this
  holds precisely for the \emph{classical} type \(\Pic^{\tau}_{Y/k} \cong \ZZ/2\ZZ\) and the \emph{supersingular}
  type \(\Pic^{\tau}_{Y/k} \cong \alpha_2\) (but not the \emph{ordinary} type \(\Pic^{\tau}_{Y/k} \cong \mu_2\)).
  The surface \(Y\) is \emph{exceptional} (in the sense of Ekedahl--Shepherd-Barron) if it is simply-connected
  and the biconductrix \(2C \subset Y\) satisfies \(h^1(\cO_{2C}) \neq 0\).
  Consequently, the \emph{non-exceptional} Enriques surfaces in characteristic \(2\) are: all ordinary surfaces
  (\(\Pic^{\tau}_{Y/k} \cong \mu_2\)) unconditionally, together with the classical (\(\ZZ/2\ZZ\)) and
  supersingular (\(\alpha_2\)) surfaces for which \(h^1(\cO_{2C}) = 0\).
\end{definition}

%-----------------------------------------------------------------
\section{Families of Enriques surfaces}
%-----------------------------------------------------------------

\begin{definition}[Picard scheme]
  \label{def:picard-scheme}
  \lean{MR4513142ThereIsNoEnriquesSurfaceOverTheIntegers.PicardScheme}
  \uses{}
  % SOURCE: [Schröer, arXiv:2004.07025], §5, lines 1208–1210 (read from references/MR4513142-enriques-integers.tex)
  % SOURCE QUOTE: "According to \cite{Artin 1969}, Theorem 7.3
  % the fppf sheafification of the naive Picard functor $A\mapsto\Pic(\foY\otimes_RA)$ is representable by a
  % group object in the category of algebraic spaces over $R$."
  \textit{Source: Schr\"oer, §5 (citing Artin 1969, Thm.~7.3).}
  Let \(R\) be a ring and \(f : \foY \to \Spec(R)\) a flat proper morphism of finite presentation.
  The \emph{Picard scheme} \(\Pic_{\foY/R}\) is the fppf sheafification of the naive Picard functor
  \(A \mapsto \Pic(\foY \otimes_R A)\), represented by a group algebraic space (in fact a scheme) over \(\Spec(R)\).
\end{definition}

\begin{definition}[Family of Enriques surfaces over a ring]
  \label{def:family-enriques-surfaces}
  \lean{MR4513142ThereIsNoEnriquesSurfaceOverTheIntegers.IsFamilyOfEnriquesSurfaces}
  \uses{def:enriques-surface}
  % SOURCE: [Schröer, arXiv:2004.07025], §5, lines 1177–1181 (read from references/MR4513142-enriques-integers.tex)
  % SOURCE QUOTE: "A \emph{family of Enriques surfaces} over an arbitrary ground ring $R$ is an algebraic space $\foY$,
  % together with a morphism  $f:\foY\ra\Spec(R)$
  % that is flat,  proper and of finite presentation such that for each geometric point $\Spec(\Omega)\ra\Spec(R)$,
  % the base-change $Y$ is an Enriques surfaces over $k=\Omega$."
  \textit{Source: Schr\"oer, §5.}
  A \emph{family of Enriques surfaces} over a ring \(R\) is an algebraic space \(\foY\) together with a morphism
  \(f : \foY \to \Spec(R)\) that is flat, proper, and of finite presentation, such that for every geometric point
  \(\Spec(\Omega) \to \Spec(R)\) the base-change \(Y = \foY \otimes_R \Omega\) is an Enriques surface over \(\Omega\).
\end{definition}

\begin{definition}[Constant Picard scheme]
  \label{def:constant-picard}
  \lean{MR4513142ThereIsNoEnriquesSurfaceOverTheIntegers.HasConstantPicardScheme}
  \uses{def:family-enriques-surfaces, def:picard-scheme, def:pic-tau, def:num-local-system}
  % SOURCE: [Schröer, arXiv:2004.07025], §5, lines 1301–1304 (read from references/MR4513142-enriques-integers.tex)
  % SOURCE QUOTE: "We say that our family of Enriques surfaces has \emph{constant local system} if
  % the local system $\Num_{\foY/R} $ is isomorphic  to
  % $(\ZZ^{\oplus 10})_R$. Likewise, we say that it has \emph{constant Picard scheme} if
  % the group space $\Pic_{\foY/R}$ is isomorphic to $(\ZZ/2\ZZ\times\ZZ^{\oplus 10})_R$."
  \textit{Source: Schr\"oer, §5.}
  A family of Enriques surfaces \(f : \foY \to \Spec(R)\) has \emph{constant Picard scheme} if
  \(\Pic_{\foY/R} \cong (\ZZ/2\ZZ \times \ZZ^{\oplus 10})_{R}\) as group schemes over \(\Spec(R)\).
  Equivalently, the local system \(\Num_{\foY/R}\) is constant (isomorphic to \((\ZZ^{\oplus 10})_R\))
  and the torsion part satisfies \(\Pic^{\tau}_{\foY/R} \cong (\ZZ/2\ZZ)_{R}\).
\end{definition}

\begin{lemma}[Simply-connectedness of \(\Spec(\ZZ)\)]
  \label{lem:spec-Z-simply-connected}
  \lean{MR4513142ThereIsNoEnriquesSurfaceOverTheIntegers.TODO.minkowskiSimplyConnected}
  \uses{}
  % SOURCE: [Schröer, arXiv:2004.07025], §5, lines 1317–1321 (read from references/MR4513142-enriques-integers.tex)
  % SOURCE QUOTE: "According to Minkowski's theorem
  % (see for example \cite{Neukirch 1999}, Theorem 2.18), the scheme $S=\Spec(\ZZ)$
  % is simply-connected"
  \textit{Source: Schr\"oer, §5, citing Neukirch 1999, Thm.~2.18.}
  The scheme \(\Spec(\ZZ)\) is simply connected: its algebraic fundamental group is trivial,
  \(\pi_1(\Spec(\ZZ), \bar{\eta}) = 1\).
  This is a consequence of the Minkowski discriminant bound.
\end{lemma}
\begin{proof}
  By the Minkowski discriminant bound, every number field \(K \neq \QQ\) has
  \(|\mathrm{disc}(K/\QQ)| > 1\), so its ring of integers is ramified over \(\ZZ\).
  In particular, the only unramified extension of \(\QQ\) is \(\QQ\) itself.
  It follows that \(\Spec(\ZZ)\) has no non-trivial connected finite \'etale covers,
  i.e.\ its algebraic fundamental group is trivial.
\end{proof}

\begin{lemma}[Vanishing of the Brauer group \(\Br(\ZZ) = 0\)]
  \label{lem:br-Z-vanishes}
  \lean{MR4513142ThereIsNoEnriquesSurfaceOverTheIntegers.TODO.brauerZVanishes}
  \uses{}
  % SOURCE: [Schröer, arXiv:2004.07025], §5, lines 1327–1333 (read from references/MR4513142-enriques-integers.tex)
  % SOURCE QUOTE: "Furthermore, we have $\Br(R)=0$. Indeed, the Hasse principle (see for example \cite{Neukirch; Schmidt; Wingberg 2008}, Theorem 8.1.17)
  % gives a short exact sequence
  % $$
  % 0\lra\Br(\QQ)\lra \bigoplus\Br(\QQ_\primid)\stackrel{\delta}{\lra} \QQ/\ZZ\lra 0,
  % $$
  % where the sum runs over all places $\primid$ including the Archimedean one. Since we have $\Br(\ZZ_p)=\Br(\FF_p)=0$ for all primes $p>0$ by Wedderburn,
  % one sees that $\Br(\ZZ)$ is contained in the Brauer group $\Br(\RR)\cap\Kernel(\delta)=0$."
  \textit{Source: Schr\"oer, §5, citing Neukirch--Schmidt--Wingberg 2008, Thm.~8.1.17.}
  The Brauer group \(\Br(\ZZ)\) vanishes.
  This follows from the global Hasse principle: the exact sequence
  \(0 \to \Br(\QQ) \to \bigoplus_v \Br(\QQ_v) \xrightarrow{\delta} \QQ/\ZZ \to 0\)
  together with \(\Br(\ZZ_p) = \Br(\FF_p) = 0\) for all primes \(p\)
  (which uses Wedderburn's theorem that every finite division ring is a field)
  implies \(\Br(\ZZ) \subseteq \Br(\RR) \cap \ker(\delta) = 0\).
\end{lemma}
\begin{proof}
  The Hasse--Brauer--Noether sequence gives an exact sequence
  \[
    0 \to \Br(\QQ) \to \bigoplus_{v} \Br(\QQ_v) \xrightarrow{\delta} \QQ/\ZZ \to 0,
  \]
  where the product runs over all places of \(\QQ\) (finite and Archimedean).
  For every finite prime \(p\), we have \(\Br(\FF_p) = 0\) by Wedderburn's theorem
  (every finite division ring is commutative).
  The Archimedean place contributes \(\Br(\RR) \cong \ZZ/2\ZZ\), but a class
  in \(\Br(\ZZ)\) (unramified at all finite places) that restricts to \(\Br(\RR)\)
  must also lie in \(\ker(\delta)\); since \(\ker(\delta) = 0\),
  we conclude \(\Br(\ZZ) = 0\).
\end{proof}

\begin{proposition}[Existence of a family of canonical coverings]
  \label{prop:canonical-covering}
  \lean{MR4513142ThereIsNoEnriquesSurfaceOverTheIntegers.canonicalCovering}
  \uses{def:family-enriques-surfaces, def:pic-tau}
  % SOURCE: [Schröer, arXiv:2004.07025], §5, lines 1245–1254 (read from references/MR4513142-enriques-integers.tex)
  % SOURCE QUOTE: "\begin{proposition}
  % \mylabel{canonical covering}
  % A family of canonical coverings $\epsilon:\foX\ra\foY$ exists
  % under any of the following three conditions:
  % \begin{enumerate}
  % \item The pullback map $H^2(S,G)\ra H^2(\foY,G_\foY)$ is injective.
  % \item The   structure morphism $f:\foY\ra S$ admits a section.
  % \item The group scheme $P\ra S$ is \'etale and $\Pic(S)=0$.
  % \end{enumerate}
  % If a family of canonical coverings exists,  it is unique up to isomorphism and twisting by  pullbacks of $G$-torsors over $S$."
  \textit{Source: Schr\"oer, §5, Prop.~5.5.}
  Let \(f : \foY \to S = \Spec(R)\) be a family of Enriques surfaces with torsion Picard scheme
  \(P = \Pic^{\tau}_{\foY/R}\) and Cartier dual \(G = \uHom(P, \GG_m)\).
  A \emph{family of canonical coverings} \(\epsilon : \foX \to \foY\) (a \(G_{\foY}\)-torsor) exists
  under any of the following three conditions:
  \begin{enumerate}
  \item the pullback map \(H^2(S, G) \to H^2(\foY, G_{\foY})\) is injective;
  \item the structure morphism \(f : \foY \to S\) admits a section;
  \item the group scheme \(P \to S\) is \'etale and \(\Pic(S) = 0\).
  \end{enumerate}
  When it exists, the family of canonical coverings is unique up to isomorphism and twisting by pullbacks of
  \(G\)-torsors over \(S\).
\end{proposition}
\begin{proof}
  Assertions (i) and uniqueness follow from the five-term exact sequence of the Leray--Serre spectral sequence
  for \(f\). A section of \(f\) provides a retraction of \(H^2(S, G) \to H^2(\foY, G_{\foY})\), giving (ii).
  For (iii), with \(P\) \'etale one takes \(\shL = \Omega^2_{\foY/R}\); the direct image \(\shN = f_*(\shL^{\otimes 2})\)
  is invertible and trivialized using \(\Pic(R) = 0\), and \(\foX = \Spec(\cO_{\foY} \oplus \shL^{\otimes -1})\)
  with the induced algebra structure is the desired covering.
\end{proof}

\begin{proposition}[Constant Picard scheme over \(\ZZ\)]
  \label{prop:constant-picard-over-Z}
  \lean{MR4513142ThereIsNoEnriquesSurfaceOverTheIntegers.constantPicardScheme}
  \uses{def:constant-picard, def:family-enriques-surfaces, lem:spec-Z-simply-connected, lem:br-Z-vanishes, prop:canonical-covering}
  % SOURCE: [Schröer, arXiv:2004.07025], §5, lines 1311–1314 (read from references/MR4513142-enriques-integers.tex)
  % SOURCE QUOTE: "\begin{proposition}
  % \mylabel{constant picard scheme}
  % If the base ring is $R=\ZZ$, then our family of Enriques surfaces $f:\foY\ra\Spec(R)$ has constant Picard scheme.
  % \end{proposition}"
  \textit{Source: Schr\"oer, §5, Prop.~5.3.}
  If \(f : \foY \to \Spec(\ZZ)\) is a family of Enriques surfaces, then \(\Pic_{\foY/\ZZ}\) is constant,
  i.e.\ \(\Pic_{\foY/\ZZ} \cong (\ZZ/2\ZZ \times \ZZ^{\oplus 10})_{\ZZ}\).
\end{proposition}

% SOURCE QUOTE PROOF: "Choose  a geometric point $\Spec(\Omega)\ra S$.
% The local system $\Num_{\foY/R}$ comes from a representation of the algebraic fundamental group
% $\pi_1(S,\Omega)$ on the stalk $\Num_{Y/k}(\Omega)=\ZZ^{\oplus 10}$. According to Minkowski's theorem
% (see for example \cite{Neukirch 1999}, Theorem 2.18), the scheme $S=\Spec(\ZZ)$
% is simply-connected, so our family of Enriques surfaces has constant local system.
%
% For the ring $R=\ZZ$, up to isomorphism the possible triples describing the group scheme $P=\Pic^\tau_{\foY/R}$
% are either $(R,1,2)$ or $(R,2,1)$. In other words, we either have $P=(\ZZ/2\ZZ)_R$ or $P=\mu_{2,R}$.
% Seeking a contradiction, we assume the latter holds.
% The Picard group $\Pic(R)$ vanishes, because the ring  $R=\ZZ$ is a principal ideal domain.
% Furthermore, we have $\Br(R)=0$. Indeed, the Hasse principle (see for example \cite{Neukirch; Schmidt; Wingberg 2008}, Theorem 8.1.17)
% gives a short exact sequence
% $$
% 0\lra\Br(\QQ)\lra \bigoplus\Br(\QQ_\primid)\stackrel{\delta}{\lra} \QQ/\ZZ\lra 0,
% $$
% where the sum runs over all places $\primid$ including the Archimedean one. Since we have $\Br(\ZZ_p)=\Br(\FF_p)=0$ for all primes $p>0$ by Wedderburn,
% one sees that $\Br(\ZZ)$ is contained in the Brauer group $\Br(\RR)\cap\Kernel(\delta)=0$.
%
% Over the localization $R'=R[1/2]$, the group scheme $\mu_{2,R}$ becomes \'etale, hence there is a canonical covering
% $\foX\otimes R'\ra\foY\otimes R'$, by Proposition \ref{canonical covering}.
% Now consider the generic fiber $X=\foX\otimes\QQ$, which is a K3 surface. In particular, the Hodge number
% $h^{2,0}(X)=h^0(\Omega_X^2)$ is non-zero.
% ...
% Recall that Fontaine showed that for each proper smooth scheme over $\QQ$ that has good reduction over
% $W(\FF_p^\sep)$ for all primes $p>0$, the Hodge numbers $h^{i,j}$ must vanish for $i\neq j$ and $i+j\leq 3$
% (\cite{Fontaine 1993}, Theorem 1 on page 44, and first remark afterwards).
% ...
% Since our K3 surface $X=\foX\otimes\QQ$ has Hodge number $h^{2,0}=1$ and   good reduction over   $W(\FF_p^\sep)$ for all $p\neq 2$, it follows that
% the base-change $X_F$ to the field of fractions $F=\Frac(A)$
% must have bad reduction over $A$.
%
% One the other hand, using Hensel's Lemma, any $\FF_2^\alg$-valued point for $\foY\ra S$ can be extended to an $A$-valued point.
% Again with Proposition \ref{canonical covering} we conclude that a family of canonical covering  $\foX\otimes A\ra\foY\otimes A$ exists,
% whence $X_F$ has good reduction over $A$, contradiction.
%
% Summing up, the Picard scheme sits in  a short exact sequence
% \begin{equation}
% \label{extension local systems}
% 0\lra (\ZZ/2\ZZ)_R\lra \Pic_{\foY/R}\stackrel{\pr}{\lra} (\ZZ^{\oplus 10})_R\lra 0.
% \end{equation}
% For each global section $l$ on the right, the preimage $T=\pr^{-1}(l)$ is a representable torsor for
% the group scheme $H=(\ZZ/2\ZZ)_R$. In particular, the structure morphism $T\ra \Spec(R)$ is finite
% \'etale, of degree two. Since $\pi_1(S,\Omega)=0$, the total space $T$ is the disjoint union of two copies
% of $S=\Spec(\ZZ)$. From this we infer that the short exact sequence \eqref{extension local systems} splits.
% In turn, our family of Enriques surfaces has constant Picard scheme."
\begin{proof}
  \uses{lem:spec-Z-simply-connected, lem:br-Z-vanishes, def:num-local-system, def:pic-tau, prop:canonical-covering}
  Since \(\Spec(\ZZ)\) is simply connected (\cref{lem:spec-Z-simply-connected}), the local system
  \(\Num_{\foY/\ZZ}\) is constant, with trivial Galois action on the stalk \(\ZZ^{\oplus 10}\).

  It remains to identify \(\Pic^{\tau}_{\foY/\ZZ}\).
  Over \(\ZZ\), the only locally free group schemes of rank two are \((\ZZ/2\ZZ)_{\ZZ}\) and \(\mu_{2,\ZZ}\).
  Assume for contradiction that \(\Pic^{\tau}_{\foY/\ZZ} = \mu_{2,\ZZ}\).
  Since \(\Pic(\ZZ) = 0\) (as \(\ZZ\) is a PID) and \(\Br(\ZZ) = 0\) (\cref{lem:br-Z-vanishes}),
  the group scheme \(\mu_{2,\ZZ}\) becomes \'etale over \(\ZZ[1/2]\), so a canonical K3-like covering
  \(\foX \otimes \ZZ[1/2] \to \foY \otimes \ZZ[1/2]\) exists by \cref{prop:canonical-covering}(iii).
  The generic fiber \(X = \foX \otimes \QQ\) is a K3 surface with \(h^{2,0}(X) = 1\),
  and it has good reduction over \(W(\FF_p^{\mathrm{sep}})\) for all odd primes \(p\).
  By Fontaine's theorem, this forces \(h^{2,0} = 0\), a contradiction.
  Therefore \(\Pic^{\tau}_{\foY/\ZZ} \cong (\ZZ/2\ZZ)_{\ZZ}\).

  The extension \(0 \to (\ZZ/2\ZZ)_{\ZZ} \to \Pic_{\foY/\ZZ} \to (\ZZ^{\oplus 10})_{\ZZ} \to 0\)
  splits because \(\pi_1(\Spec(\ZZ)) = 1\) forces every degree-two finite étale torsor over \(\Spec(\ZZ)\)
  to be trivial.
  Hence \(\Pic_{\foY/\ZZ} \cong (\ZZ/2\ZZ \times \ZZ^{\oplus 10})_{\ZZ}\).
\end{proof}

\begin{theorem}[No exceptional Enriques surface over \(\ZZ/4\ZZ\)]
  \label{thm:no-exceptional-over-Z4}
  \lean{MR4513142ThereIsNoEnriquesSurfaceOverTheIntegers.noExceptionalOverZ4}
  \uses{def:exceptional-enriques, def:family-enriques-surfaces}
  % SOURCE: [Schröer, arXiv:2004.07025], §5, lines 1199–1203 (read from references/MR4513142-enriques-integers.tex)
  % SOURCE QUOTE: "\begin{theorem}
  % \mylabel{no exceptional}
  % There is no family of Enriques surfaces over the ring $R=\ZZ/4\ZZ$ whose
  % geometric fiber is exceptional.
  % \end{theorem}"
  \textit{Source: Schr\"oer, §5, Thm.~5.2 (citing Schr\"oer 2021b, Thm.~7.2).}
  There is no family of Enriques surfaces over \(R = \ZZ/4\ZZ\) whose geometric fiber is exceptional.
\end{theorem}

\begin{proof}
  This is cited from \cite{Schroeer 2021b}, Theorem 7.2, and is not reproved in the present formalization.
  It is stated here as a formal axiom (or admitted sorry) to make the dependency graph explicit.
\end{proof}

\begin{theorem}[Family with constant Picard scheme has a genus-one fibration]
  \label{thm:family-with-constant-pic}
  \lean{MR4513142ThereIsNoEnriquesSurfaceOverTheIntegers.hasGenusOneFibration}
  \uses{def:constant-picard, def:family-enriques-surfaces, def:picard-scheme, def:pic-tau, prop:canonical-covering, thm:projective-contractions}
  % SOURCE: [Schröer, arXiv:2004.07025], §5, lines 1394–1416 (read from references/MR4513142-enriques-integers.tex)
  % SOURCE QUOTE: "\begin{theorem}
  % \mylabel{family with constant pic}
  % Suppose that the family of Enriques surfaces  $f:\foY\ra\Spec(R)$ has constant Picard scheme.
  % Assume furthermore  that  $\Pic(R)$  and  $\Br(R)$ vanish.
  % Then there is at least one canonical covering $\epsilon:\foX\ra \foY$, and
  % the set of isomorphism classes  of such coverings
  % is a principal homogeneous space for the  group
  % $R^\times/R^{\times 2}$. Furthermore, the induced Enriques surfaces $Y$ and $\bY$ enjoy the following properties:
  % ...
  % \item
  % There is at least one genus-one fibration $\varphi:Y\ra\PP^1$.
  % \end{enumerate}
  % \end{theorem}"
  \textit{Source: Schr\"oer, §5, Thm.~5.4.}
  Suppose \(f : \foY \to \Spec(R)\) is a family of Enriques surfaces with constant Picard scheme,
  and that \(\Pic(R) = \Br(R) = 0\).
  Let \(k\) be a residue field of \(R\) and \(Y = \foY \otimes_R k\).
  Then:
  \begin{enumerate}
  \item The dualizing sheaf \(\omega_Y\) has order two in the Picard group.
  \item Every numerical class in \(\Num_{Y/k}(\bar{k})\) comes from an invertible sheaf on \(Y\),
        unique up to twisting by \(\omega_Y\).
  \item Every \((-2)\)-curve on \(\bar{Y} = Y \otimes_k \bar{k}\) descends to a \((-2)\)-curve on \(Y\),
        isomorphic to \(\PP^1_k\).
  \item Every genus-one fibration on \(\bar{Y}\) descends to a genus-one fibration \(\varphi : Y \to \PP^1_k\)
        with exactly two tame multiple fibers over \(k\)-rational points.
  \item There exists at least one genus-one fibration \(\varphi : Y \to \PP^1_k\).
  \end{enumerate}
\end{theorem}

% SOURCE QUOTE PROOF: "\proof
% By assumption we have $\Pic^\tau_{\foY/R}=(\ZZ/2\ZZ)_R$.
% Write $S=\Spec(R)$. The Kummer sequence $0\ra \mu_2\ra\GG_m\stackrel{2}\ra\GG_m\ra 0$ gives an exact sequence
% $$
% \Pic(S)\stackrel{2}{\lra}\Pic(S)\lra H^2(S,\mu_{2})\lra \Br(S)\stackrel{2}{\lra}\Br(S).
% $$
% Our assumptions ensure that the term in the middle vanishes.
% According to Proposition \ref{canonical covering},   there is at least one canonical covering $\epsilon:\foX\ra\foY$.
% ...
% It remains to verify (iv).
% Let $\bY\ra\PP^1_{\bar{k}}$ be a genus-one fibration.
% According to Theorem \ref{projective contractions}, it is the base-change of a fibration $Y\ra B$, where $B$ is a twisted form
% of the projective line.  Consider the invertible sheaf  $\bar{\shL}=\O_\bY(\bar{F})$, where $\bar{F}\subset\bY$
% is a half-fiber. It arises from some invertible sheaf $\shL$ on $Y$.   Using $h^0(\bar{\shL})=1$ we infer that there is a unique curve
% $F\subset Y$ inducing $\bar{F}\subset \bY$. ...
% Summing up, we have shown that each half-fiber on $\bY$ induces a half-fiber on $Y$ that lies over a rational point.
% Since there are exactly two half-fibers on $\bY$, both of which are tame, the same holds for $Y$."
\begin{proof}
  \uses{def:constant-picard, prop:canonical-covering, thm:projective-contractions}
  Since the Picard scheme is constant, we have \(\Pic^{\tau}_{\foY/R} \cong (\ZZ/2\ZZ)_R\).
  The Kummer sequence for \(\mu_2\) and the vanishing of \(\Pic(R)\) and \(\Br(R)\) imply that
  \(H^2(\Spec(R), \mu_2) = 0\), so a canonical covering \(\foX \to \foY\) exists (\cref{prop:canonical-covering}).

  For conclusion (v): Since \(Y\) has constant Picard scheme, every numerical class over \(\bar{k}\) descends to \(Y\),
  and by the structure theory of Enriques surfaces there exists at least one genus-one fibration on \(\bar{Y}\).
  By \cref{thm:projective-contractions} (projective contractions descend when the local system is constant),
  this fibration descends to a fibration \(\varphi : Y \to \PP^1_k\),
  with multiple fibers descending to tame fibers over rational points.
\end{proof}

\begin{theorem}[No Enriques surface over the integers]
  \label{thm:no-enriques-over-integers}
  \lean{MR4513142ThereIsNoEnriquesSurfaceOverTheIntegers.noEnriquesSurfaceOverIntegers}
  \uses{prop:constant-picard-over-Z, thm:no-exceptional-over-Z4, thm:no-non-exceptional-over-F2}
  % SOURCE: [Schröer, arXiv:2004.07025], §5, lines 1188–1192 (read from references/MR4513142-enriques-integers.tex)
  % SOURCE QUOTE: "\begin{theorem}
  % \mylabel{no enriques over integers}
  % The fiber category $\shM_\Enr(\ZZ)$ is empty. In other words,
  % there is no family of Enriques surfaces over the ring $R=\ZZ$.
  % \end{theorem}"
  \textit{Source: Schr\"oer, §5, Thm.~5.1 (proved in §15).}
  The fiber category \(\shM_{\Enr}(\ZZ)\) is empty.
  In other words, there is no family of Enriques surfaces over the ring \(R = \ZZ\).
\end{theorem}

% SOURCE QUOTE PROOF: "We are finally ready to prove Theorem \ref{no enriques over integers},  the main result of this paper.
% The task is to show that the fiber category $\shM_\Enr(\ZZ)$ for the stack of Enriques surfaces is empty.
% Seeking a contradiction, we assume that that there is a family of Enriques surfaces
% $f:\foY\ra \Spec(\ZZ)$. According to Proposition \ref{constant picard scheme}, the Picard scheme $\Pic_{\foY/\ZZ}$ is constant.
% So the closed fiber $Y=\foY\otimes\FF_2$ also has constant Picard scheme.
% It is  non-exceptional, by  \cite{Schroeer 2021b}, Theorem 7.2, because this Enriques surface
%  comes with a deformation over $W_2(\FF_2)=\ZZ/4\ZZ$.
% Therefore it suffices to establish:
% [Theorem: no non-exceptional Enriques surface over FF_2 with constant Picard scheme.]"
\begin{proof}
  \uses{prop:constant-picard-over-Z, thm:no-exceptional-over-Z4, thm:no-non-exceptional-over-F2}
  Suppose for contradiction that there exists a family \(f : \foY \to \Spec(\ZZ)\) of Enriques surfaces.

  By \cref{prop:constant-picard-over-Z}, the Picard scheme \(\Pic_{\foY/\ZZ}\) is constant.
  In particular, the closed fiber \(Y = \foY \otimes_{\ZZ} \FF_2\) has constant Picard scheme.

  The surface \(Y\) is non-exceptional: it admits a deformation over \(W_2(\FF_2) = \ZZ/4\ZZ\)
  (namely the base-change \(\foY \otimes_{\ZZ} \ZZ/4\ZZ\)), and by \cref{thm:no-exceptional-over-Z4}
  no exceptional Enriques surface admits such a deformation.

  But \cref{thm:no-non-exceptional-over-F2} asserts that no Enriques surface over \(\FF_2\)
  is simultaneously non-exceptional and has constant Picard scheme. Contradiction.
\end{proof}

%-----------------------------------------------------------------
\section{Classification of certain jacobian fibrations}
%-----------------------------------------------------------------

\begin{theorem}[Classification of Weierstrass equations]
  \label{thm:classification-weierstrass}
  \lean{MR4513142ThereIsNoEnriquesSurfaceOverTheIntegers.classificationWeierstrassEqns}
  \uses{def:constant-picard, def:picard-scheme}
  % SOURCE: [Schröer, arXiv:2004.07025], Introduction, lines 523–529 (read from references/MR4513142-enriques-integers.tex)
  % SOURCE QUOTE: "Up to isomorphisms, there are exactly eleven Weierstra\ss{}
  % equations $y^2+a_1xy + \ldots =x^3+a_2x^2+\ldots$ with coefficients $a_i\in \FF_2[t]$
  % that define a geometrically rational elliptic surface $\phi:J\ra\PP^1$
  % with constant Picard scheme $\Pic_{J/\FF_2}$  having at most one rational point $a\in \PP^1$ where $J_a$ is semistable or supersingular."
  \textit{Source: Schr\"oer, Introduction / §10 (Thms.~10.1--10.2).}
  Up to isomorphism, there are exactly eleven Weierstra\ss{} equations
  \(y^2 + a_1 x y + \cdots = x^3 + a_2 x^2 + \cdots\) with coefficients \(a_i \in \FF_2[t]\)
  that define a geometrically rational elliptic surface \(\phi : J \to \PP^1\)
  over \(\FF_2\) with constant Picard scheme \(\Pic_{J/\FF_2}\),
  having at most one rational point \(a \in \PP^1\) where \(J_a\) is semistable or supersingular.
\end{theorem}
\begin{proof}
  The proof is an explicit enumeration given in Section~10 of the paper (Theorems~10.1 and~10.2 there).
  One writes down all Weierstrass equations \(y^2 + a_1 x y + a_3 y = x^3 + a_2 x^2 + a_4 x + a_6\)
  with \(a_i \in \FF_2[t]\) that define geometrically rational elliptic surfaces over \(\FF_2\)
  (i.e.\ minimal Weierstrass fibrations \(J \to \PP^1\) with geometric genus zero base and section).
  The constancy of \(\Pic_{J/\FF_2}\) translates into arithmetic conditions on the coefficients.
  The constraint that at most one rational point \(a \in \PP^1\) is semistable or supersingular
  (a condition coming from the reduction to Enriques surfaces) further limits the possibilities.
  A direct computer-aided enumeration, checking all valid combinations of
  Kodaira symbols for the fibers at \(t = 0, 1, \infty\) and verifying discriminant conditions,
  yields exactly eleven isomorphism classes.
  These are recorded explicitly in the two classification theorems.
\end{proof}

%-----------------------------------------------------------------
\section{Proof of the main result}
%-----------------------------------------------------------------

\begin{theorem}[No non-exceptional Enriques surface over \(\FF_2\) with constant Picard scheme]
  \label{thm:no-non-exceptional-over-F2}
  \lean{MR4513142ThereIsNoEnriquesSurfaceOverTheIntegers.noNonExceptionalOverF2}
  \uses{def:enriques-surface, def:exceptional-enriques, def:constant-picard,
        thm:family-with-constant-pic, thm:classification-weierstrass}
  % SOURCE: [Schröer, arXiv:2004.07025], §15, lines 4088–4092 (read from references/MR4513142-enriques-integers.tex)
  % SOURCE QUOTE: "\begin{theorem}
  % \mylabel{no non-exceptional}
  % There is no  Enriques surface $Y$ over   $k=\FF_2$  that is  non-exceptional and  whose   Picard scheme
  % is constant.
  % \end{theorem}"
  \textit{Source: Schr\"oer, §15, Thm.~15.1.}
  There is no Enriques surface \(Y\) over the prime field \(k = \FF_2\) that is simultaneously
  non-exceptional and has constant Picard scheme \(\Pic_{Y/k} \cong (\ZZ^{\oplus 10} \oplus \ZZ/2\ZZ)_k\).
\end{theorem}

% SOURCE QUOTE PROOF: "Suppose such a $Y$ exists.
% By Theorem \ref{family with constant pic} there is a genus-one fibration $\varphi:Y\ra\PP^1$.
% For each geometric point $\ba:\Spec(\Omega)\ra \PP^1$ lying over a non-rational closed point  $a\in\PP^1$, the geometric fiber $Y_\ba$
% is  irreducible, according to Proposition \ref{constant num over finite field},  thus the  Kodaira symbol is $\I_0, \I_1$ or $\II$.
% Moreover, it must be simple by Theorem \ref{family with constant pic}.
% The possible configurations of fibers over the rational points $t=0,1,\infty$ are tabulated
% in Proposition \ref{key reduction}.
%  In Sections \ref{Elimination I4*} and \ref{Elimination III*} we saw that the Kodaira symbols
% $\I_4^*$ and $\III^*$ actually do not occur.
% For the classical Enriques surface $\bY=Y\otimes_kk^\alg$   this means that for every genus-one fibration, the configuration
% of degenerate fibers is either $\I_1^*+\I_4$ or $\I_2^*+\III+\III$. Furthermore, in the former case the $\I_1^*$-fiber is multiple,
% whereas in the latter case two of the three fibers are multiple.
%  We showed in Proposition \ref{only two configurations} that
% such an  Enriques surface does not exist. This gives the desired contradiction."
\begin{proof}
  \uses{thm:family-with-constant-pic, thm:classification-weierstrass}
  Suppose such a \(Y\) exists.
  By \cref{thm:family-with-constant-pic} (applied with \(R = \FF_2\), where \(\Pic(\FF_2) = \Br(\FF_2) = 0\)),
  there exists a genus-one fibration \(\varphi : Y \to \PP^1\).
  For geometric points over non-rational closed points \(a \in \PP^1\), the geometric fiber is irreducible
  (by a point-counting argument using constancy of \(\Num\)), so its Kodaira symbol is \(\mathrm{I}_0\), \(\mathrm{I}_1\), or \(\mathrm{II}\);
  it must also be simple.
  The possible configurations of degenerate fibers over the three rational points \(t = 0, 1, \infty\)
  are classified using \cref{thm:classification-weierstrass} (passage to the Jacobian fibration \(J\))
  together with a theorem of Liu--Lorenzini--Raynaud on the relation between an elliptic fibration and its Jacobian.
  The Kodaira symbols \(\mathrm{I}_4^*\) and \(\mathrm{III}^*\) are eliminated in Sections~12--13 of the paper.
  The remaining configurations reduce to \(\mathrm{I}_1^* + \mathrm{I}_4\) or \(\mathrm{I}_2^* + \mathrm{III} + \mathrm{III}\)
  on the base-change to \(\bar{\FF}_2\).
  An extensive combinatorial analysis of pairs of genus-one fibrations (Section~14) shows that no Enriques surface
  admits these configurations. Contradiction.
\end{proof}
